\backhead{FURTHER EXERCISES}{khred}
\begin{exlist}
\item Fresh manure has $k=0.05$ d$^{-1}$ at 20 $^\circ$C. Find its half-life at 20 $^\circ$C and at 10 $^\circ$C ($Q_{10}=2$).\exans{13.9 d; 27.7 d}
\item Clover residues (C:N 15) of 3 t DM\,ha$^{-1}$ (42 per cent C) decompose fully with $e=0.4$ and $CN_m=8$. Find the net N released.\exans{21 kg N\,ha$^{-1}$}
\item A herbicide has $K_{oc}=100$ L\,kg$^{-1}$; the soil has 1.5 per cent organic carbon, $\rho_b=1.4$ kg\,L$^{-1}$, $\theta=0.25$ and drains 250 mm\,yr$^{-1}$. Find $K_d$, $R$ and the time to travel 1 m.\exans{1.5 L\,kg$^{-1}$; 9.4; 9.4 yr}
\item A field drains 180 mm at 15 mg NO$_3$-N\,L$^{-1}$. Find the N leached and the mean nitrate concentration in mg NO$_3$\,L$^{-1}$.\exans{27 kg N\,ha$^{-1}$; 66 mg NO$_3$\,L$^{-1}$}
\item Nitrification follows first-order kinetics with $k=0.1$ d$^{-1}$. How long does it take to nitrify 90 per cent of an ammonium application?\exans{23 d}
\end{exlist}
\backhead{ESSAY QUESTIONS}{khred}
{\small\itshape Answer each question in about one page, with equations, a sketch and at least one number from this chapter.}\par\vspace{3pt}
\begin{exlist}
\item Explain the pools of organic matter in Daisy and the meaning of their decay rates.
\item Discuss mineralisation and immobilisation and the role of the C:N ratio of residues.
\item Explain the processes of nitrate loss (leaching, denitrification) and how irrigation affects each.
\item Describe the transport of a pesticide through the soil and the factors that decide whether it reaches the groundwater.
\item Discuss nitrous oxide emission from irrigated soils and how it could be reduced.
\item Explain how the nitrogen balance of a field is closed and checked.
\end{exlist}
\dykband{EXERCISES IN SWAFL}
{\small Open SWAFL (desktop or online) with the benchmark project. The chapter of the \href{https://github.com/halilunaziru73-creator/Coupled-Water-Fluxes-Pore-Scale-Dynamics-Precision-Irrigation-Mulching-Feedback-Loops/raw/main/Software/v1.1.0/SWAFL_Desktop/docs/manual_v1.1.0/SWAFL_Manual_v1.1.0.pdf}{SWAFL manual} that explains each task is given in brackets; record your values in a table and compare them with what is expected (\emph{Check}).}\par\vspace{3pt}
\begin{exlist}
\item {\color{khblue}[Manual, Chapter 22]} Open \emph{Nitrogen budget} and read the leaching of the wettest and the driest year.\par\noindent{\small\itshape Check: leaching is largest in the wettest year.}
\item {\color{khblue}[Manual, Chapter 28]} Design a Daisy experiment with three nitrogen rates and run it; compare the leaching.\par\noindent{\small\itshape Check: leaching rises with the rate.}
\item {\color{khblue}[Manual, Chapter 30]} Open the results of the chemical scenario and read the pesticide leached.\par\noindent{\small\itshape Check: most is degraded near the surface.}
\item {\color{khblue}[Manual, Chapter 25]} Click every box of the nitrogen part of the process map and list the processes.\par\noindent{\small\itshape Check: mineralisation, nitrification, denitrification, leaching, uptake.}
\end{exlist}
\backhead{OBJECTIVE TYPE QUESTIONS}{khred}
\begin{exlist}
\item The critical C:N ratio above which residues immobilise nitrogen is about\opt{5}{20}{80}{200}
\item The retardation factor of a non-sorbed solute is\opt{0}{1}{2}{$\rho_b/\theta$}
\item Denitrification is favoured by\opt{a dry soil}{a well-aerated soil}{a wet soil}{a frozen soil}
\item 50 mg NO$_3$\,L$^{-1}$ equals about\opt{3.5}{11.3}{25}{50} mg N\,L$^{-1}$
\item The half-life of a first-order process with $k=0.1$ d$^{-1}$ is\opt{0.69 d}{6.9 d}{10 d}{69 d}
\end{exlist}
\objkey{1.\ (b),\quad 2.\ (b),\quad 3.\ (c),\quad 4.\ (b),\quad 5.\ (b)}
